\chapter{Conclusion}
\label{sec:conclusion}

This thesis evaluated prompt decomposition strategies for telematics-based
fleet monitoring using two local open-source LLMs:
\texttt{llama3.1:8b} and \texttt{deepseek-r1:8b}. The central task was graph
recovery from Tesla EV telemetry. Each model-strategy output was parsed into
directed feature edges, converted into an adjacency matrix, and compared with
a six-edge physics-informed Tesla EV ground-truth graph.

The contribution is a controlled, artifact-traceable evaluation rather than a
claim of autonomous causal discovery. Five prompt strategies were compared:
zero-shot prompting, no-instruction control, policy-first structured
prompting, least-to-most prompting, and skeleton-of-thought prompting.
Performance was reported with precision, recall, F1 score, SHD, FDR, Pearson
correlation, Hamming distance, and Frobenius norm.

The strongest observed result is \texttt{deepseek-r1:8b} with
skeleton-of-thought prompting. This configuration achieves precision 0.4444,
recall 0.6667, F1 score 0.5333, SHD 7, FDR 0.5556, Pearson correlation
0.4303, Hamming distance 0.1944, and Frobenius norm 2.6458. It recovers four
of the six ground-truth edges but also predicts five false positive edges.
This is the clearest evidence that decomposition can help in the evaluated
Tesla EV setting, but it is not reliable enough to be treated as validated
causal knowledge.

The broader results show model-specific trade-offs. Llama has higher average
F1 across strategies, but it also produces a higher average false discovery
rate and more structural error. DeepSeek produces the best single graph, but
three of its validated strategy cells contain empty outputs and zero parsed
edges. Strategy effects are also mixed. Skeleton-of-thought is the strongest
decomposition strategy in the supplied results; least-to-most prompting does
not improve the evaluated graph recovery. Policy-first instruction structuring
is not sufficient on its own.

The consensus graph reinforces the need for evaluation. It captures correct
recurring edges such as \texttt{electricenergyremaining -> electricremaining},
but it also preserves recurring false edges such as
\texttt{speed\_kmh -> acceleration\_kmh2} and mileage-distance reversals.
Agreement across prompts and models can therefore reflect shared mistakes,
not only shared knowledge.

The final claim is precise. Prompt decomposition can improve
LLM-assisted telematics graph recovery in this controlled Tesla EV setting,
especially when the prompt structure matches the feature-system structure.
However, the evaluated LLMs remain responsive to prompt wording, output format,
parser behavior, and model-specific generation patterns. The thesis does not
establish true causal discovery, repeated-run self-consistency, or general
reliability across vehicle datasets.

For fleet-monitoring research, the durable contribution is an auditable
evaluation discipline for LLM-assisted telemetry analytics. A responsible
workflow should define the feature set, constrain the model to produce a
parseable graph, convert the proposal into an adjacency matrix, compare it
with explicit domain knowledge, and inspect false positives and direction
errors before operational use. Under this discipline, LLMs are best understood
as structured graph proposal systems: useful because they can surface
candidate relationships, limited because those proposals require validation,
and accountable only when human reviewers retain final responsibility for the
interpretation and use of the graph.
